
\documentclass{report}

\input{preamble}
\input{macros}
\input{letterfonts}

\title{\Huge{Economia Aziendale}\\Sistemi Informativi Aziendali}
\author{\huge{Caberlotto Giovanni}}
\date{}

\begin{document}

\maketitle
\newpage% or \cleardoublepage
% \pdfbookmark[<level>]{<title>}{<dest>}
\pdfbookmark[section]{\contentsname}{toc}
\tableofcontents
\pagebreak

\chapter{Società di Capitali}
\dfn{Società di Capitali}{Le società di capitali si dividono in: \begin{itemize}
    \item Società per azioni s.p.a.
    \item Società in accomandita per azioni s.a.p.a.
    \item Società a responsabilità limitata s.r.l.
    \item Società a responsabilità limitata semplificata s.r.l.s.
\end{itemize} Si tratta di organizzazioni di persone e mezzi per l’esercizio in comune di un’attività produttiva, dotate di piena autonomia patrimoniale: quindi, soltanto la società con il suo patrimonio risponde delle obbligazioni sociali. Il socio, pertanto, ha una responsabilità limitata al capitale conferito, non assumendo alcuna responsabilità personale, neanche sussidiaria, per le obbligazioni sociali (tranne i casi previsti dalla legge)}

\section{Costituzione}

    

\qs{Costituzione S.r.l.}{In data 1/4 viene fondata la società Alfa S.r.l., capitale sociale 2.000.000, composto da 1.000.000€ in denaro e 1.000.000€ da un macchinario, il cui valore è confermato dalla relazione giurata. Il denaro viene conferito tramite assegni, versati il giorno stesso dagli amministratori sul c/c bancario. In data 1/7 gli amministratori richiamano il restante 75\% in denaro tramite lettera raccomandata, che viene versato il 15/7}

\begin{table}[h]
    \centering
    \caption{Libro Giornale}
    \begin{tabular}{|l|l|l|l|}
    \hline
    Data & Descrizione & Dare & Avere \\
    \hline
    1/04 & Soci c/sottoscrizione & 2.000.000,00 & \\
         & Capitale sociale & & 2.000.000,00 \\
    \hline
    1/04 & Impianti e macchinari & 1.000.000,00 & \\
         & Assegni & 250.000,00 & \\
         & Soci c/sottoscrizione & & 1.250.000,00 \\
    \hline
    1/04 & Banca c/c & 250.000,00 & \\
         & Assegni & & 250.000,00 \\
    \hline
    1/07 & Soci c/versamenti richiami & 750.000,00 & \\
         & Azionisti c/sottoscrizione & & 750.000,00 \\
    \hline
    15/07 & Banca x c/c & 750.000,00 & \\
          & Azionisti c/versamenti richiami & & 750.000,00 \\
    \hline
    \end{tabular}
\end{table}

    
\qs{Costituzione S.p.a.}{In data 1/4 viene fondata la società Alfa S.p.A., capitale sociale 2.000.000 suddiviso in 100.000 azioni da 20€. Il 50\% è composto da conferimenti in denaro e l'altro 50\% da un macchinario, il cui valore di 1.000.000€ è confermato dalla relazione giurata. Il denaro viene inviato tramite assegni, versati il giorno stesso dagli amministratori sul c/c vincolato. In data 13/4 avviene la registrazione presso il registro delle imprese e vengono svincolati i fondi, fruttiferi al tasso del 1,5\%. In data 1/7 gli amministratori richiamano il restante 75\% in denaro tramite lettera raccomandata, che viene versato il 15/7}

\begin{table}[h]
    \centering
    \caption{Libro Giornale}
    \begin{tabular}{|l|l|l|l|}
    \hline
    Data & Descrizione & Dare & Avere \\
    \hline
    1/04 & Soci c/sottoscrizione & 2.000.000,00 & \\
         & Capitale sociale & & 2.000.000,00 \\
    \hline
    1/04 & Impianti e macchinari & 1.000.000,00 & \\
         & Assegni & 250.000,00 & \\
         & Soci c/sottoscrizione & & 1.250.000,00 \\
    \hline
    1/04 & Banca c/c & 250.000,00 & \\
         & Assegni & & 250.000,00 \\
    \hline
    1/07 & Soci c/versamenti richiami & 750.000,00 & \\
         & Azionisti c/sottoscrizione & & 750.000,00 \\
    \hline
    15/07 & Banca x c/c & 750.000,00 & \\
          & Azionisti c/versamenti richiami & & 750.000,00 \\
    \hline
    \end{tabular}
\end{table}

\dfn{Definizione Parcella}{Nota delle competenze, delle spese fatte o di altri compensi spettanti di diritto, che viene presentata al cliente da un libero professionista}

\qs{Parcella notaio}{
Si riceve in data 16/4 la parcella del notaio per la redazione dell'atto costitutivo, che reca 4.000€ di onorario, 750€ di spese documentate, IVA 22\% e ritenuta fiscale 20\%. Il collegio sindacale approva la patrimonializzazione del costo. La parcella viene pagata il 18/4 tramite bonifico.
}
\begin{table}[h]
    \centering
    \caption{Libro Giornale}
    \begin{tabular}{|l|l|l|l|}
    \hline
    Data & Descrizione & Dare & Avere \\
    \hline
    16/04 & Costi d'impianto & 4750,00 & \\
         & Iva ns/credito & 880,00 &  \\
         & Debiti vs/fornitori & & 5630,00 \\
    \hline
    16/04 & Debiti vs/fornitori & 5630,00 & \\
         & Debiti per ritenuta da versare & 800,00 & \\
         & Banca c/c & & 4830,00 \\
    \hline
    \end{tabular}
\end{table}
\nt{Quando si parla di parcelle bisogna distingure due casistiche: \begin{itemize}
    \item Parcella del notaio;
    \item Parcella del commercialista.
\end{itemize}
nella seconda andrà aggiunta in fattura all'imponibile iva il costo "contributo cassa di prevvidenza" del 4\%}
\qs{Parcella Commercialista}{Il giorno 10/1 si riceve la parcella n. 9 del commercialista, con onorario di 3000€ + contributo di previdenza 4\% + IVA. Sono anche addebitate spese documentate per 400€. Viene pagata immediatamente }

\begin{table}[h]
    \centering
    \caption{Registrazione Parcella Commercialista}
    \begin{tabular}{|l|l|}
        \hline
         Onorario & 3.000,00  \\
         Contributo cassa di prevvidenza (4\%)& 120 \\
         Imponibile & 3.120,00 \\
         Iva (22\%) & 686,40 \\
         Spese documentate & 400 \\
         \hline
         Totale Parcella & 4.206,40 \\ 
         \hline
         Ritenuta (20\%) su onorario & 600,00 \\
         \hline
         Totale da pagare & 3.606,40 \\
         \hline
    \end{tabular}
\end{table}

\begin{table}[h]
    \centering
    \caption{Libro Giornale}
    \begin{tabular}{|l|l|l|l|}
    \hline
    Data & Descrizione & Dare & Avere \\
    \hline
    16/04 & Costi d'impianto & 3.520,00 & \\
         & Iva ns/credito & 686,40 &  \\
         & Debiti vs/fornitori & & 4.206,40 \\
    \hline
    16/04 & Debiti vs/fornitori & 4.206,40 & \\
         & Debiti per ritenuta da versare &  & 600,00 \\
         & Banca c/c & & 3606,40 \\
    \hline
    \end{tabular}
\end{table}

\section{Riparto utili e copertura perdite}
\qs{Riparto utili}{In data 12/4 viene approvato il progetto di riparto dell'utile dell'anno n - che ammonta a 146.000€ - e che prevede l'accantonamento a riserva legale e il 3\% a riserva statutaria. Sono inoltre presenti utili a novo per 14890€. Il capitale sociale è composto da 20.000 azioni dal valore di 30€, di cui 15.000 possedute da soggetti lordisti. Il pagamento viene eseguito il 28/4}

\nt{
Il calcolo del dividendo è dato dall'utile da distribuire diviso dalle quantità di azioni, per il calcolo dell'utile da distribuire agli azionisti basterà invece moltiplicare il risultato del rapporto tra utile da distribuire e quantità di azioni per il numero totale di azioni:
$ \\ Dividendo \  =  \ \frac{Utile \ da \ distribuire}{Numero \ di \ azioni} \\
Dividendo  \ Totale \ = \ Dividendo \ * \ Numero \ di \ azioni \\ $ 
}

\begin{table}[h]
    \centering
    \caption{Prospetto riparto utile}
    \begin{tabular}{|l|l|l|}
    \hline
         Utile d'esercizio & & 146000  \\
    \hline
         Riserva Legale & 7300 & \\ 
         Riserva Statutaria & 4380 & \\
    \hline
        A Riserve & & -11680 \\
        Utili a Nuovo & & +14890 \\
    \hline
        Utile da distribuire & & 149210 \\
        Dividendi & & 149200 \\
    \hline
        Utile a nuovo & & 10 \\
    \hline
    \end{tabular}
\end{table}

\begin{table}[h]
    \centering
    \caption{Libro Giornale}
    \begin{tabular}{|l|l|l|l|}
    \hline
    Data & Descrizione & Dare & Avere \\
    \hline
    12/4 & Utile d'esercizio &146.000,00 & \\
         &  Utile a nuovo & 14.890,00 &  \\
         &  Riserva Legale & & 7.300,00\\
         &  Riserva Statutaria & & 4.380,00\\
         & Azionisti c/dividendi & & 149.200,00\\
         & Utile a nuovo & & 10,00 \\
    \hline
    16/04 & Azionisti c/dividendi & 149.200,00 & \\
         & Debiti per ritenuta da versare &  & 9.700,00 \\
         & Banca c/c & & 139.500,00 \\
    \hline
    \end{tabular}
\end{table}

\nt{La ritenuta fiscale del 26\%viene applicata solamente sui nettisti e viene calcolata tramite la seguente formula:
$ \\ Ritenuta \ unitaria \ = \ Dividendo \ * 26\% \\
Ritenuta \ Totale \ = \ Ritenuta \ unitaria \ * \ num.nettisti$}

\qs{Perdita rimandata all'anno successivo}{Il 31/12/n si rileva una perdita di 23.450€, che viene rimandata all'anno successivo. L'utile di esercizio dell'anno n+1 risulta di 59.850€ e viene usato per coprire le perdite portate a nuovo.}

\begin{table}[h]
    \centering
    \caption{Libro Giornale}
    \begin{tabular}{|l|l|l|l|}
    \hline
    Data & Descrizione & Dare & Avere \\
    \hline
    16/04 & Perdite a nuovo & 23.450,00  & \\
         & Perdite d'esercizio &  &  23.450,00\\
    \hline
    13/04/n+1 & Utile d'esercizio & 59.850,00& \\
         & Riserva Legale &  & 2.992,50 \\
         & Perdite a nuovo & & 23.450,00 \\
         & Azionisti c/dividendi & & 33.407,50 \\
    \hline
    \end{tabular}
\end{table}

\qs{Copertura perdite tramite riserve}{Per coprire le perdite di 49.590 vengono utilizzate la riserva straordinaria, pari a 30.000€, e la riserva legale, pari a 129.000€}
\newpage
\begin{table}[h]
    \centering
    \caption{Libro Giornale}
    \begin{tabular}{|l|l|l|l|}
    \hline
    Data & Descrizione & Dare & Avere \\
    \hline
    16/04 & Riserva Legale &  19.590,00 & \\
          & Riserva statutaria & 30.000,00 & \\
          & Perdite d'esercizio &  & 49.590,00  \\
    \hline
    \end{tabular}
\end{table}

\qs{Copertura perdite tramite riduzione capitale sociale}{Per coprire le perdite di 467.980€, si approva la riduzione del capitale sociale per 300.000€ e il reintegro di capitale per il restante}

\begin{table}[h]
    \centering
    \caption{Libro Giornale}
    \begin{tabular}{|l|l|l|l|}
    \hline
    Data & Descrizione & Dare & Avere \\
    \hline
    16/04 & Capitale Sociale &  300.000,00 & \\
          & Perdita d'esercizio &  & 300.000,00 \\
    \hline
          & Banca c/c & 167.980,00 & \\
          & Azionisti c/dividendi & & 167.980,00 \\
    \hline
    \end{tabular}
\end{table}
\nt{La copurterua della perdita può essere coperta anche tramite il fondo utile a nuovo}
\nt{La copertura della perdita può prevedere anche l'uso di tutti questi casi contemporaneamente}

\section{Aumento capitale sociale}


\end{document}
